\section{Análisis del Verification Prompt}

\subsection{El papel y las instrucciones del Verifier}

El System Prompt de Verification también es Markdown estructurado y contiene:

\begin{enumerate}
    \item \textbf{Interpretación de un papel}: eres un matemático excepcional, un revisor con el rigor de nivel IMO
    \item \textbf{Indicación central}: \textbf{eres solo el Verifier, no el Solver}: no intentes corregir errores ni llenar gap, solo verifica step by step
    \item \textbf{Clasificación de problemas}:
    \begin{itemize}
        \item \textbf{Critical Error}: error lógico grave (por ejemplo, de $A > B, C > D$ deducir $A-C > B-D$)
        \item \textbf{Gap}: paso omitido, ruptura lógica
    \end{itemize}
    \item \textbf{Formato de salida}: Summary (con Verdict: valid/invalid) + List of Findings (bug report) + Detailed Verification Log
\end{enumerate}

\begin{warningbox}{El Verifier no debe intentar resolver el problema}
El Verification Prompt exige explícitamente que el Verifier no intente corregir errores ni llenar gap, y que solo verifique. Esta \textbf{separación de papeles} es clave en el diseño del Workflow: el Generator y el Verifier cumplen cada uno su función.
\end{warningbox}

\subsection{Resumen de la sección}

Mediante una definición estricta de papeles y un formato de salida estructurado, el Verification Prompt garantiza que el Verifier solo verifique y no resuelva; su Verdict y su Bug Report resultantes impulsan la Correction posterior y la decisión de terminar el ciclo.

